\documentclass[11pt,fleqn]{article} 
\usepackage[margin=0.8in, head=0.8in]{geometry} 
\usepackage{amsmath, amssymb, amsthm}
\usepackage{fancyhdr} 
\usepackage{palatino, url, multicol}
\usepackage{graphicx, pgfplots} 
\usepackage[all]{xy}
\usepackage{polynom,tabularx} 
%\usepackage{pdfsync} %% I don't know why this messes up tabular column widths
\usepackage{enumerate}
\usepackage{framed}
\usepackage{setspace}
\usepackage{array}
\usepackage{pgf,tikz}
\usepackage{mathrsfs}

\usepackage[parfill]{parskip}
\usetikzlibrary{arrows}

\usetikzlibrary{calc}

\pgfplotsset{compat=1.6}

\pgfplotsset{soldot/.style={color=blue,only marks,mark=*}} \pgfplotsset{holdot/.style={color=blue,fill=white,only marks,mark=*}}

\renewcommand{\headrulewidth}{0pt}
\newcommand{\blank}[1]{\rule{#1}{0.75pt}}
\newcommand{\bc}{\begin{center}}
\newcommand{\ec}{\end{center}}
\newcommand{\be}{\begin{enumerate}}
\newcommand{\ee}{\end{enumerate}}

\def\ds{\displaystyle}

\renewcommand{\d}{\displaystyle}

\newcommand{\ans}[1][2]{ \ \rule{#1 in}{.5 pt} \ }


\pagestyle{fancy} 
%\lfoot{Uses a calculator}
\rfoot{4-10}

\begin{document}

\vspace*{-0.7in}

\begin{center}
  \Large\sc{Review: Derivative and Integration Rules}
\end{center}

The left column is for differentiation rules. The right column is for the corresponding integration rule \emph{if such a rule exists}.\\

\doublespacing

\begin{tabularx}{\textwidth}{l || l} 
Differentiation Rules & Integration Rules\\ \hline \hline & \\
\hspace*{-.3in}(a) \: $\displaystyle{\frac{d}{dx}\left(\tan(x)\right)=}$ \hspace{2.5in} \quad & \\ &\\ \hline & \\
\hspace*{-.3in}(b) \:  $\displaystyle{ \frac{d}{dx}\left(k\cdot g(x)\right)=}$& \\ 
 $k$ is a constant\\ &\\ \hline & \\
\hspace*{-.3in}(c) \:  $\displaystyle{\frac{d}{dx}\left(e^x\right)=}$ &\\ &\\ \hline & \\
\hspace*{-.3in}(d) \:  $\displaystyle{ \frac{d}{dx}\left(x^k\right)=}$ &\\ &\\ \hline & \\ 
\hspace*{-.3in}(e) \:  $\displaystyle{ \frac{d}{dx}\left(\ln(x)\right)=}$ &\\ &\\ \hline & \\
\hspace*{-.3in}(f) \:  $\displaystyle{ \frac{d}{dx}\left(x^k\right)=}$ &\\ 
$k \not = -1$ &\\ \hline & \\
\hspace*{-.3in}(g) \:  $\displaystyle{ \frac{d}{dx}\left(\sin(x)\right)=}$&\\ &\\ \hline & \\
\hspace*{-.3in}(h) \: $\displaystyle{ \frac{d}{dx}\left(\cos(x)\right)=}$ &\\ &\\ \hline & \\
\hspace*{-.3in}(i) \:  $\displaystyle{ \frac{d}{dx}\left(\sec(x)\right)=}$ &\\ 
 \end{tabularx}
 
 \newpage
 
\begin{tabularx}{\textwidth}{l || l} 
Differentiation Rules & Integration Rules\\ \hline \hline & \\
\hspace*{-.3in}(j) \: $\displaystyle{ \frac{d}{dx}\left(\arcsin(x)\right)=}$ \hspace{2.5in} \quad & \\ &\\ \hline & \\
\hspace*{-.3in}(k) \:  $\displaystyle{\frac{d}{dx}\left(c\right)=}$& \\ 
 $c$ is a constant\\  \hline & \\
\hspace*{-.3in}(l) \:  $\displaystyle{\frac{d}{dx}\left(\arctan(x)\right)=}$ &\\ &\\ \hline & \\
\hspace*{-.3in}(m) \:  $\displaystyle{ \frac{d}{dx}\left(f(x)\cdot g(x)\right)=}$ &\\ &\\ \hline & \\ 
\hspace*{-.3in}(n) \:  $\displaystyle{ \frac{d}{dx}\left(\frac{f(x)}{g(x)}\right)=}$ &\\ &\\ \hline & \\
\hspace*{-.3in}(o) \:  $\displaystyle{\frac{d}{dx}\left(f(g(x))\right)=}$ &\\  &\\ \hline & \\
\hspace*{-.3in}(p) \:  $\displaystyle{ \frac{d}{dx}\left({f(x)}+{g(x)}\right)=}$&\\ &\\ \hline & \\
\hspace*{-.3in}(q) \: $\displaystyle{\frac{d}{dx}\left(\csc(x)\right)=}$ &\\ &\\ \hline & \\
\hspace*{-.3in}(r) \:  $\displaystyle{ \frac{d}{dx}\left(\cot(x) \right)=}$ &\\ &\\ \hline & \\
\hspace*{-.3in}(s) \:  $\displaystyle{ \frac{d}{dx}\left(2^x \right)=}$ &\\ 
 \end{tabularx}
 
 \end{document}
 


  && \\
&&\\
\hline
  & $\displaystyle{ \frac{d}{dx}\left({f(x)}+{g(x)}\right)=}$\\
&where $k$ is a constant&\\
\hline
 $\displaystyle{\frac{d}{dx}\left(\csc(x)\right)=}$ &$\displaystyle{ \frac{d}{dx}\left(\cot(x) \right)=}$&$\displaystyle{ \frac{d}{dx}\left(|x| \right)=}$\\
&&\\

 \end{tabular}

\singlespacing
\item Write the equivalent integral formula \emph{where possible.}

\end{enumerate}
\end{document}
